%% ============================================================
%% teampage.tex — 제1부 도입: 팀 구성 및 업무 분장 (독립 장 아님)
%% 제1부 표지 다음 한 페이지를 비우고 홀수 페이지에서 표로 표시된다.
%% ============================================================
\thispagestyle{fancy}
\addcontentsline{toc}{chapter}{팀 구성 및 업무 분장}

\begin{center}
  {\sffamily\bfseries\LARGE\color{deepblue}팀 구성 및 업무 분장}
\end{center}
\vspace{1.4em}

\noindent 교수진과 제품개발학기 참여 학생의 역할 분담은 다음과 같다. 서비스 애플리케이션,
시뮬레이션, 펌웨어, UI/UX, 실물 교구 디자인을 나누어 진행한다.

\vspace{1.6em}
\noindent{\sffamily\bfseries\large\color{aresorange}교수진}
\begin{center}
\renewcommand{\arraystretch}{1.5}
\begin{tabularx}{\linewidth}{@{}>{\sffamily\bfseries}l L{0.74\linewidth}@{}}
\toprule
\sffamily 강영민 & 하드웨어 제어 및 시뮬레이션 SW 개발총괄 \\
\sffamily 서미라 & 제품 디자인 지도 (실물 교구 및 온라인 서비스) \\
\sffamily 신\,\,원 & 제품개발학기 총괄 진행 및 제품 수준의 서비스 관리 \\
\bottomrule
\end{tabularx}
\end{center}

\vspace{1.2em}
\noindent{\sffamily\bfseries\large\color{aresorange}학생}
\begin{center}
\renewcommand{\arraystretch}{1.5}
\begin{tabularx}{\linewidth}{@{}>{\sffamily\bfseries}l L{0.74\linewidth}@{}}
\toprule
\sffamily 김지훈 & 시뮬레이션 모듈 고도화 및 디지털 트윈 시스템 구현 \\
\sffamily 이민혁 & 시뮬레이션 모듈 고도화 및 디지털 트윈 시스템 구현 \\
\sffamily 이주현 & 펌웨어 개발 및 블록코딩 서비스 연계 \\
\sffamily 이성빈 & 제품 서비스 UI/UX 고도화 \\
\sffamily 이재훈 & 실물 교구 디자인 기획 및 시작품 제작 \\
\bottomrule
\end{tabularx}
\end{center}
